% !TeX program = pdflatex
%==============================================================================
%  CauldronHook — Uniswap v4 Hook Routing Allowlist: Reviewer Brief
%
%  Companion to (and deliberately SEPARATE from) the independent security audit
%  in MagicFrens_Independent_Audit_2026-09.tex. That document is a neutral
%  assessment; this one is a submission brief. Keeping them apart is the point —
%  an audit that doubles as advocacy is worth less to a reviewer than an audit
%  that does not.
%
%  Build:  tectonic MagicFrens_Uniswap_Hook_Allowlist_Brief.tex
%     or:  pdflatex MagicFrens_Uniswap_Hook_Allowlist_Brief.tex  (x2)
%==============================================================================
\documentclass[11pt,a4paper]{article}

\usepackage[T1]{fontenc}
\usepackage[utf8]{inputenc}
\usepackage[margin=2.4cm]{geometry}
\usepackage{booktabs}
\usepackage{longtable}
\usepackage{array}
\usepackage{xcolor}
\usepackage{colortbl}
\usepackage{fancyhdr}
\usepackage{titlesec}
\usepackage{enumitem}
\usepackage{listings}
\usepackage[hidelinks,bookmarks=true]{hyperref}
\usepackage{amssymb}
\usepackage{tcolorbox}
\tcbuselibrary{breakable,skins}

\definecolor{ink}{HTML}{16181D}
\definecolor{slate}{HTML}{5B6472}
\definecolor{rule}{HTML}{D5D9E0}
\definecolor{good}{HTML}{15803D}
\definecolor{warn}{HTML}{A16207}
\definecolor{bad}{HTML}{9B1C1C}
\definecolor{uni}{HTML}{FF007A}
\definecolor{panel}{HTML}{F6F7F9}

\titleformat{\section}{\Large\bfseries\color{ink}\raggedright}{\thesection}{0.7em}{}
\titleformat{\subsection}{\normalsize\bfseries\color{ink}\raggedright}{\thesubsection}{0.6em}{}
\setlist{leftmargin=1.35em,itemsep=2pt,topsep=3pt}

\pagestyle{fancy}\fancyhf{}
\fancyhead[L]{\small\color{slate}CauldronHook --- v4 Routing Allowlist Brief}
\fancyhead[R]{\small\color{slate}\thepage}
\renewcommand{\headrulewidth}{0.4pt}

\lstdefinestyle{sol}{
  basicstyle=\ttfamily\scriptsize, breaklines=true, columns=fullflexible,
  keywordstyle=\color{uni}, commentstyle=\color{slate}\itshape,
  numbers=none, frame=leftline, framerule=1.5pt, rulecolor=\color{rule},
  xleftmargin=0.9em, showstringspaces=false, language=Java,
}
\lstset{style=sol}

\newcommand{\code}[1]{\texttt{\small #1}}
\newcommand{\Yes}{\textcolor{good}{\textbf{Yes}}}
\newcommand{\No}{\textcolor{bad}{\textbf{No}}}

\newtcolorbox{keybox}[1]{
  breakable, enhanced, colback=panel, colframe=uni, boxrule=1pt,
  left=7pt,right=7pt,top=6pt,bottom=6pt, arc=1pt,
  title={#1}, fonttitle=\bfseries\small, coltitle=white
}
\newtcolorbox{plainbox}[1]{
  breakable, enhanced, colback=panel, colframe=rule, boxrule=0.7pt,
  left=7pt,right=7pt,top=6pt,bottom=6pt, arc=1pt,
  title={#1}, fonttitle=\bfseries\small\color{ink}, coltitle=ink
}

\sloppy
\begin{document}

%==============================================================================
\begin{center}
{\LARGE\bfseries\color{ink} CauldronHook\par}
\vspace{2pt}
{\large\color{slate} Uniswap v4 Hook Routing Allowlist --- Reviewer Brief\par}
\vspace{6pt}
{\small\color{slate} Magic Internet Frens \quad$\cdot$\quad \href{https://mifrens.xyz}{mifrens.xyz}\par}
\end{center}
\vspace{6pt}

\begin{keybox}{Why this submission exists}
\code{CauldronHook} declares \code{beforeSwapReturnDelta} and
\code{afterSwapReturnDelta}. Both are structurally required by the hook's single
economic mechanism --- an \textbf{ETH-denominated} swap fee --- for reasons that
follow from v4's specified/unspecified currency semantics rather than from a
design preference. Section~\ref{sec:why} shows the derivation.

We do \textbf{not} deploy at \code{0x91...}, and we do \textbf{not} target major
pairs: every pool this hook serves is \code{ETH / <our own iteration token>},
created by our own registry.
\end{keybox}

\vspace{4pt}

\begin{center}\small
\begin{tabular}{@{}lp{0.62\linewidth}@{}}
\toprule
\multicolumn{2}{@{}l}{\textbf{Answers to the stated criteria}} \\
\midrule
Upgradeable? & \No{} --- no proxy, no delegatecall in the hook, immutable bytecode (\S\ref{sec:upgrade}) \\
Requires custom data inputs? & \No{} --- \code{hookData} is optional on every path; empty is a first-class case (\S\ref{sec:hookdata}) \\
Audited? & \Yes{} --- independent from-scratch review, 340 tests incl.\ live-v4 fork tests (\S\ref{sec:audit}) \\
Differentiated? & \Yes{} --- \S\ref{sec:diff} \\
Deployment address & \code{0x31ee...50cc} --- not \code{0x91} \\
Pairs & ETH / own token only --- never a major pair \\
\bottomrule
\end{tabular}
\end{center}

%==============================================================================
\section{What the hook is}
\label{sec:diff}

The Cauldron is an autonomous, self-funding token machine. One \emph{iteration}
token trades at a time; when its rolling 24-hour volume falls below a threshold,
governance selects the next one, protocol-owned liquidity is recovered, and
holders migrate 1:1. The cycle is designed to run indefinitely, and the protocol
is funded entirely by its own swap fees.

A break-glass custody path does exist and we would rather state it than have it
found: \code{emergencyWithdrawLP} can remove a generation's protocol-owned
liquidity to an emergency admin. That admin is \textbf{immutable} (fixed at
construction, no setter), the action must be announced on-chain and is delayed, an
independent guardian can veto it in the window, and arming it \emph{forces the
redemption exit open} so holders can leave at the floor before anything moves. It
reaches protocol-owned liquidity only --- nothing in the system can mint, freeze,
or move a token or NFT held in a user's wallet, and none of it is reachable from
the hook. Full detail in \S\ref{sec:upgrade}.

Four things are genuinely v4-native rather than ported from a v2/v3 design:

\begin{itemize}
\item \textbf{The fee engine is the hook.} There is no LP fee
(\code{POOL\_FEE = 0}); 100\% of trading revenue is the hook's ETH-denominated
return-delta fee, routed per-swap to an NFT-holder dividend, a per-collection
token floor, the iteration's proposer, and the reserve that funds the next
launch.
\item \textbf{Hook-native perpetuals.} Leveraged longs and shorts whose opens,
closes and liquidations are \emph{real pool swaps}, with no external oracle.
Liquidation marks come from the engine's own on-chain TWAP ring. Underwater
positions are swept opportunistically from \code{afterSwap}, so any swap on any
interface can trigger a liquidation and earn the keeper reward.
\item \textbf{Progressive in-swap seeding.} Launch liquidity streams into the
pool over a window using \emph{core} \code{modifyLiquidity} from inside
\code{afterSwap} --- keeperless, and un-frontrunnable because the deployed
fraction is a pure function of elapsed time. Anti-snipe by book \emph{shape}
rather than by tax alone.
\item \textbf{Commit--reveal NFT forging in \code{afterSwap}.} A direct Uniswap
or aggregator swap --- with no router and no \code{hookData} --- forges NFTs
in-swap. This is a deliberate first-class path, not a fallback (\S\ref{sec:hookdata}).
\end{itemize}

%==============================================================================
\section{Permissions, and why the delta flags are required}
\label{sec:why}

\subsection{Declared permissions}

\begin{lstlisting}
afterInitialize:       true    // adopt only pools our registry created
beforeSwap:            true    // charge the ETH fee on buys
afterSwap:             true    // charge the ETH fee on sells
beforeSwapReturnDelta: true
afterSwapReturnDelta:  true
// everything else false: no initialize/donate/liquidity callbacks
\end{lstlisting}

\noindent Encoded flag bits: \code{0x10cc} = \code{01000011001100}. The deployed
address \code{0x31ee3436d96533e91874cda5dfc70cb521f750cc} carries exactly those
low bits and no others --- verifiable with
\code{uint160(hook) \& Hooks.ALL\_HOOK\_MASK}. We hold no liquidity-callback or
donate permissions, so the hook cannot interpose on \code{modifyLiquidity} or
\code{donate} at all.

\subsection{The derivation}

The hook's economics require the fee to be denominated in \textbf{ETH on every
trade}, because the token side of each pool is by construction a
soon-to-be-retired asset --- accruing protocol revenue in it would strand that
revenue at the next rebirth.

In v4, \code{afterSwap}'s return delta can only act on the \textbf{unspecified}
currency. For our pools (\code{currency0} is native ETH):

\begin{center}\small
\begin{tabular}{@{}llll@{}}
\toprule
\textbf{Quadrant} & \textbf{Specified} & \textbf{Unspecified} & \textbf{\code{afterSwap} could charge in} \\
\midrule
\rowcolor{panel} exact-in BUY \, (\code{zeroForOne}, \code{amt<0}) & ETH (in) & token (out) & \textcolor{bad}{token} \\
exact-out BUY (\code{zeroForOne}, \code{amt>0}) & token (out) & ETH (in) & \textcolor{good}{ETH} \\
exact-in SELL (\code{!zeroForOne}, \code{amt<0}) & token (in) & ETH (out) & \textcolor{good}{ETH} \\
\rowcolor{panel} exact-out SELL (\code{!zeroForOne}, \code{amt>0}) & ETH (out) & token (in) & \textcolor{bad}{token} \\
\bottomrule
\end{tabular}
\end{center}

The \textbf{exact-input buy} --- the dominant trade, and what the Uniswap
interface and every aggregator emit --- has the \emph{token} as its unspecified
leg. Charging it in ETH therefore requires skimming the specified input in
\code{beforeSwap}, i.e.\ \code{beforeSwapReturnDelta}. There is no alternative
mechanism.

\code{afterSwapReturnDelta} charges the sell side. Without it the hook cannot
take value from a swap at all: a \code{poolManager.take()} with no offsetting
return delta leaves an unsettled negative delta and the unlock reverts
\code{CurrencyNotSettled}.

\begin{plainbox}{We evaluated the delta-free alternative and rejected it on merit}
A dynamic LP fee (\code{DYNAMIC\_FEE\_FLAG} plus a \code{beforeSwap} override)
would need no delta flag, and our LP is protocol-owned, so the fee would land in
our own positions. We did not choose it, for reasons we think are the right ones
rather than convenient ones:

\begin{enumerate}[leftmargin=1.6em]
\item \textbf{LP fees accrue in the input currency}
(\code{Pool.sol}: \code{zeroForOne ? feeGrowthGlobal0 : feeGrowthGlobal1}), so
every \emph{sell} would accrue the retiring token --- exactly the outcome the
ETH-only fee exists to prevent.
\item \textbf{It would introduce a keeper.} Our positions sit in the periphery
\code{PositionManager}, which opens its own \code{unlock} and therefore cannot be
harvested inside a swap. Per-swap routing would become a batch harvest.
\item \textbf{It would be JIT-extractable.} We hold no liquidity-callback
permissions, so LPing into our pools is permissionless. Under an LP fee a
searcher could JIT a concentrated position around a large swap and capture the
fee --- including the launch surtax, inverting the mechanic it exists for. Under
the delta model this is structurally impossible, because the hook takes the fee
directly regardless of who provides liquidity.
\end{enumerate}
The delta flags are therefore load-bearing, and the alternative would make the
hook \emph{less} safe, not more.
\end{plainbox}

%==============================================================================
\section{Integration behaviour a router should know}
\label{sec:integration}

We would rather you learn these from us than find them in review.

\subsection{Disclosed: exact-output sells currently revert}

\code{\_beforeSwap} reverts \code{ExactOutSellUnsupported} for the
(\code{amountSpecified > 0}, \code{!zeroForOne}) quadrant. The reason is the table
in \S\ref{sec:why}: the unspecified leg there is the token, so v4's return-delta
mechanism physically cannot take an ETH fee. Rather than silently serve that
quadrant \emph{for free} --- which was the prior behaviour, and a measured fee
bypass --- we refuse it.

\begin{plainbox}{Remediation, in progress}
We accept that reverting a whole quadrant is a poor router citizen, and it is the
one item in this brief we consider a genuine integration defect rather than a
trade-off. Because the specified currency in that quadrant \emph{is} ETH, the fee
can be taken via \code{beforeSwapReturnDelta} on the specified side --- which we
already hold, so the fix requires \textbf{no additional permission} and no
address re-mine. We are implementing it and will confirm before any mainnet
routing goes live. Every other quadrant is served today.
\end{plainbox}

\subsection{Fee bounds}

\begin{center}\small
\begin{tabular}{@{}llp{0.46\linewidth}@{}}
\toprule
\textbf{Component} & \textbf{Value} & \textbf{Notes} \\
\midrule
Base hook fee & 300 bps (3\%) default & Governance-tunable, hard-capped at \code{MAX\_TAX\_BPS = 1000} (10\%) \\
Launch surtax & up to 9{,}600 bps & \textbf{Time-decaying anti-snipe}; linear decay to \textbf{0} across \code{snipeWindowBlocks} (default 30 blocks). Routed 100\% to NFT holders \\
Absolute clamp & 9{,}900 bps & \code{MAX\_TOTAL\_FEE\_BPS} --- a swap always retains $\geq$1\% so it can never return zero \\
LP fee tier & \textbf{0} & Traders pay the hook fee only; there is no stacked LP fee \\
\bottomrule
\end{tabular}
\end{center}

\noindent We flag the surtax explicitly because a hook that can take 99\% of a
swap deserves a second look. It exists only in the opening minutes of a freshly
summoned pool, decays deterministically to zero, is bounded by an immutable
\code{constant} ceiling, and is paid to token holders rather than to us. Steady
state is 3\%.

\subsection{\code{afterSwap} does optional work --- none of it can break a swap}

\code{afterSwap} may trigger a nested buyback swap, a seeder liquidity
placement, a perp-liquidation sweep, and an NFT mint. Every one of those is
invoked with a \textbf{reserved gas budget and an ignored result}:

\begin{lstlisting}
uint256 g = gasleft();
if (g > LIQ_GAS_RESERVE + LIQ_GAS_MIN) {
    perpEngine.call{gas: g - LIQ_GAS_RESERVE}(...);   // result ignored
}
\end{lstlisting}

\noindent so an optional step that runs out of gas or reverts is a silent no-op,
and the parent swap always completes its fee take and returns. Reserves are held
for the liquidation sweep, the gacha step, the buyback, and the seeder poke
independently. Nested swaps are guarded by transient flags
(\code{\_inSelfBuy}, \code{\_inRelaunchClose}) so they never re-enter the fee
logic or recurse.

\subsection{Delta settlement}

Four contracts may settle inside one unlock (hook, registry library, perp engine,
seeder). Each settles its own deltas in its own frame. Two details the audit
specifically verified:

\begin{itemize}
\item The buyback settles the \textbf{realised} debit read from the returned
delta, not the intended amount --- so a binding price limit cannot leave a
positive un-taken delta and revert the user's swap.
\item The perp engine switches to an \textbf{in-place} swap when already inside
the lock, rather than re-entering \code{unlock}.
\end{itemize}

\subsection{Pool adoption is closed}

v4 lets anyone initialise a pool naming any hook, so \code{\_afterInitialize} is
an untrusted entrypoint. We adopt a pool only if \code{sender == registry}
\textbf{and} \code{currency0} is native ETH. A foreign pool naming this hook is
\emph{declined rather than reverted} --- so we cannot be used to grief pool
creation --- and every value-moving path is gated on the adoption flag, meaning
such a pool receives no fee logic, no accounting and no spending.

%==============================================================================
\section{Not upgradeable}
\label{sec:upgrade}

The hook is a plain contract. No proxy, no \code{delegatecall} in the hook, no
admin that can replace its code. Three things could superficially read as
upgradeability, so we describe them precisely:

\begin{enumerate}
\item \textbf{Policy modules} (\code{IDeathChecker}, \code{ISurtaxPolicy},
\code{IOddsPolicy}, \code{ICurvePolicy}, \code{IFeeRouter}). Each is a
\emph{view} that returns a number. It never custodies or moves funds, its output
is clamped to the same immutable ceilings as the built-in rule, and a reverting
module falls back to the built-in rule. \code{IFeeRouter} returns \emph{amounts}
only --- the hook performs the sends, and a split that does not sum exactly to
the fee is discarded in favour of the built-in split. This is how economic
parameters are tuned without an upgradeable money contract.
\item \textbf{Controller pointer.} \code{setRegistry} is one-shot. Replacing it
requires \code{proposeRegistryOverride} $\rightarrow$ a \code{REGISTRY\_SWAP\_DELAY}
of \textbf{7 days}, declared \code{constant} so the owner cannot shorten their own
notice period $\rightarrow$ \code{executeRegistryOverride}, which \emph{flushes}
the accrued fee reserve to the \emph{outgoing} controller first, so a swap cannot
capture ETH the previous controller accrued.
\item \textbf{Ownership.} The hook is owned by an OpenZeppelin
\code{TimelockController} --- an audited standard, chosen specifically to avoid
an upgradeable proxy on the money path. Every parameter change is
schedule\,$\rightarrow$\,wait\,$\rightarrow$\,execute.
\end{enumerate}

\subsection{The break-glass path, stated in full}

Governance \textbf{can} tune fees and risk parameters, swap a policy module, pause
a redemption flow, and --- through an announced, delayed, guardian-vetoable
action --- \textbf{remove protocol-owned liquidity} via
\code{emergencyWithdrawLP} / \code{emergencySweep} on the registry. We list this
explicitly because it is a genuine custody power and any reviewer reading the
source will find it.

Four things bound it:
\begin{itemize}
\item The emergency admin is \textbf{immutable} --- set in the constructor with
no setter, so it cannot be re-pointed at a fresh wallet after the fact. It is the
governance timelock.
\item Actions \textbf{must be announced on-chain} (\code{armEmergency}) before they
can execute, and are then delayed. Our audit found that this was previously
conditional on a non-zero delay --- a zero-delay deployment skipped the arm, which
silently disabled the forced-open exit below. It is now unconditional, so the
guarantee holds at any configuration (audit finding F-19, fixed).
\item An independent \textbf{guardian can veto} during that window. The guardian
can only ever cancel --- never propose, never move funds.
\item \textbf{Arming forces the exit open.} The redemption circuit-breaker is
overridden the instant an action is armed, so it cannot be used to trap holders
ahead of a custody move.
\end{itemize}

\noindent Governance \textbf{cannot} mint tokens, freeze balances, pause
transfers, touch any token or NFT held in a user's wallet, or shorten its own
announced delays. \textbf{None of this is reachable from the hook} --- these are
registry functions over protocol-owned liquidity, and they cannot affect swap
execution, settlement, or any routed trade.

%==============================================================================
\section{No required custom data inputs}
\label{sec:hookdata}

\code{hookData} is \textbf{optional on every code path}, and an empty
\code{hookData} is a designed-for case rather than a degraded one:

\begin{itemize}
\item \textbf{Fee tier} --- \code{\_taxedPlayer} falls back to \code{sender} when
\code{hookData} is absent.
\item \textbf{NFT credit} --- falls back to \code{tx.origin}, precisely so that a
buyer arriving through the Uniswap interface or an aggregator still earns credit
rather than having it credited to the router.
\item \textbf{Fee exemption} --- requires \code{isOpener[sender]}, an explicit
allowlist. A direct swapper controls \code{hookData} freely, so an arbitrary
caller cannot tag an exempt address and dodge the fee.
\end{itemize}

\noindent The in-swap NFT forging path exists \emph{specifically} for
router-less swaps. A swap with empty \code{hookData} is fully served: it is
charged correctly, it accrues credit, it can forge an NFT, and it can trigger a
perp liquidation.

%==============================================================================
\section{Audit and testing}
\label{sec:audit}

\begin{itemize}
\item \textbf{Independent audit} --- a from-scratch review of all $\sim$9{,}200
lines, conducted against the intended L2 deployment target, deriving its findings
solely from source. Companion document:
\code{MagicFrens\_Independent\_Audit\_2026-09.pdf}. No critical findings; every
Medium and Low finding either fixed with a regression test that fails against the
pre-fix bytecode, or disclosed with remediation.
\item \textbf{340 tests, 66 suites}, all passing: unit, pure fuzz,
handler-driven invariants (128{,}000 calls per invariant), and \textbf{fork tests
that drive live Uniswap v4} rather than a mock.
\item \textbf{Verified structural properties}, argued from cited code in the
audit's invariant register: fixed non-mintable supply; migration exactly 1:1 or
it reverts; every v4 delta settled on all four settlement paths; pool adoption
closed; and byte-identical storage layouts across the registry's delegatecall
pair (53 entries, verified mechanically).
\end{itemize}

\subsection{Disclosed limitations}

We would rather state these than have them found:

\begin{center}\small
\begin{tabular}{@{}p{0.30\linewidth}p{0.62\linewidth}@{}}
\toprule
\textbf{Item} & \textbf{Status} \\
\midrule
Exact-output sells revert & Integration defect. Fix identified, needs no new permission (\S\ref{sec:integration}). \\
Reserve ceiling & Our redemption reserve is a single-sided band; if the token appreciates past it, redemption and migration pause until the next iteration. Protocol-internal --- it does not affect swap routing or settlement. \\
\code{blockhash}-derived NFT rarity & Sound on Ethereum; weaker on our L2 target, where \code{blockhash} is documented as insecure. Being verified against the target chain before that feature is enabled there. Does not affect swaps. \\
EIP-170 headroom & The hook sits 88 bytes below the limit. A process constraint, not a user-facing risk. \\
\bottomrule
\end{tabular}
\end{center}

%==============================================================================
\section{One operational question}

Our design summons a \textbf{new pool at every rebirth} --- same hook, same
\code{ETH / <iteration token>} shape, new \code{currency1} and therefore a new
\code{PoolId}. The hook itself is permanent and immutable.

We have assumed allowlisting attaches to the \textbf{hook address}, which would
make this submission a one-time review. If it instead attaches per-pool, we would
appreciate knowing, since it materially affects how we operate --- and we would
rather align with your process now than discover the mismatch after launch.

%==============================================================================
\section{Artifacts}

\begin{center}\footnotesize
\begin{tabular}{@{}l@{\hspace{1.4em}}l@{}}
\toprule
Hook (current deployment) & \code{0x31ee3436d96533e91874cda5dfc70cb521f750cc} \\
Permission bits & \code{0x10cc} --- afterInitialize, before/afterSwap, both swap deltas \\
Pair shape & \code{ETH / <iteration token>}, registry-created only \\
Independent audit & \code{audit/MagicFrens\_Independent\_Audit\_2026-09.pdf} \\
Hook source & \code{CauldronHook.sol} \quad (repo: \code{contracts/solidity/}) \\
Test suite & \code{FOUNDRY\_PROFILE=cauldron forge test} \\
Public docs & \href{https://mifrens.xyz/\#/docs}{mifrens.xyz/\#/docs} \, ($\cdot$ machine-readable at \code{/llms.txt}) \\
\bottomrule
\end{tabular}
\end{center}

\vspace{10pt}
\begin{plainbox}{In one paragraph}
\code{CauldronHook} uses both swap delta flags because its only revenue
mechanism is an ETH-denominated fee, and v4's unspecified-currency semantics make
\code{beforeSwapReturnDelta} the sole way to charge an exact-input buy in ETH. It
is not upgradeable, requires no custom data inputs, serves only pools its own
registry created against native ETH, never touches a major pair, holds no
liquidity or donate permissions, and every optional action it performs in
\code{afterSwap} is gas-bounded and result-ignored so it cannot revert or starve a
routed swap. It has been independently audited and is covered by 339 tests
including live-v4 fork tests. The one known integration defect --- a reverting
exact-output-sell quadrant --- is disclosed above with a fix that requires no
additional permission, as is the registry's announced, delayed and
guardian-vetoable break-glass over protocol-owned liquidity, which is not
reachable from the hook and cannot affect a routed swap.
\end{plainbox}

\end{document}
